Conditional inference trees choose a split variable by a permutation test of
association before searching for a threshold, a separation designed to reduce
the preference of CART-style splitting for features with many cutpoints. We ask
how conditional inference trees (CIT) and forests (CIF), whose splits pass a
Stage~A feature test and a Stage~B threshold test each under a Bonferroni
correction, compare with established feature rankers, and what that correction
costs and changes. We benchmark them against 15 classification and 16
regression rankers, among them partykit's ctree and cforest, random forests
(RF), recursive feature elimination with random forests (RF-RFE), ExtraTrees,
and gradient boosting, on 31 real datasets, scoring each ranking by the
held-out accuracy of fixed learners fitted to its top-$k$ features. CIF ranks
tied 3rd of 17 in classification, behind RF-RFE and XGBoost and level with RF;
2nd of 18 in regression, behind ExtraTrees. On the smaller panels where every
method completed it ranks tied 6th and tied 5th, the regression ordering is
descriptive (Friedman $p=0.072$), and single trees rank near the bottom. Three
answers are negative: on synthetic designs with known informative features CIF
sits in the lower half on recovery (8th of 15 and 11th of 16), it is weakest at
small $k$, and a random forest ranked by out-of-bag permutation importance does
no better than CIF or the impurity-ranked forest. At the minimum permutation
budget the correction costs on the order of $m^2/\alpha$ permutations per node in the
feature test over $m$ features and $K^2/\alpha$ in the per-threshold rule over
$K$ candidate thresholds; the threshold term dominates on most datasets and the
feature term on the widest, where one tree on 5,000 features takes 29 minutes.
A max-type threshold test, an opt-in alternative to the default per-threshold
rule, reaches the nominal level with about $999K$ statistic evaluations, has
realized size closer to nominal, and loses 0.006 in power when the two rules
are compared at equal realized size (95\% interval $-0.011$ to $-0.002$); under
it CIF ranks 5th of 17 and 4th of 18. The adaptive stopping rule has exact null
size 0.0488 at a fixed budget of 999 permutations for a continuous null. For a
practitioner, CIF gives a real-data ranking level with random forests in
classification, and the max-type rule lowers its fitting cost at nodes with
many candidate thresholds.
